\documentclass[aspectratio=169,10pt]{beamer}
\usepackage[utf8]{inputenc}
\usepackage[spanish]{babel}
\usepackage{tikz}
\usepackage{booktabs}
\usepackage{xcolor}
\usetikzlibrary{shapes.geometric,arrows.meta,positioning,backgrounds,calc}

\usetheme{Madrid}
\setbeamercolor{palette primary}{bg=blue!65!black,fg=white}
\setbeamercolor{palette secondary}{bg=blue!45!black,fg=white}
\setbeamercolor{structure}{fg=blue!65!black}
\setbeamercolor{block title}{bg=blue!12,fg=blue!65!black}
\setbeamercolor{alerted text}{fg=orange!85!black}
\setbeamertemplate{navigation symbols}{}

\newcommand{\src}[1]{{\tiny\color{gray}Fuente: #1}}

\title[Tablas de Par\'ametro en Power BI]%
      {Por qu\'e las Tablas de Par\'ametro\\No Tienen Peso y S\'i Tienen Sentido}
\subtitle{VertiPaq, SWITCH y dise\~no din\'amico de matrices}
\author{Departamento de Planeaci\'on --- Antioquia}
\date{Mayo 2026}

\begin{document}

% =============================================================
\begin{frame}
\titlepage
\end{frame}

% =============================================================
\begin{frame}{La pregunta}

\begin{alertblock}{Lo que se pregunta}
``Si una tabla no tiene relaci\'on con los datos reales,
\textbf{?`para qu\'e sirve?
?`No agrega peso innecesario al modelo?}''
\end{alertblock}

\bigskip
Para responder hay que entender tres cosas:

\begin{enumerate}
  \item C\'omo Power BI guarda los datos en memoria (\textbf{motor VertiPaq})
  \item Qu\'e \textbf{problema real} resuelve la tabla de par\'ametro
  \item C\'omo funciona la f\'ormula \texttt{SWITCH + SELECTEDVALUE}
\end{enumerate}

\bigskip
\src{Russo \& Ferrari, \textit{The Definitive Guide to DAX}, 2nd ed., Microsoft Press, 2019}

\end{frame}

% =============================================================
\begin{frame}{C\'omo guarda los datos Power BI: motor VertiPaq}

Power BI \textbf{no guarda filas}. Guarda \textbf{columnas} comprimidas de forma independiente.

\bigskip
\begin{center}
\begin{tikzpicture}[
  col/.style={draw=blue!40, fill=blue!8, rounded corners=2pt,
              minimum width=1.3cm, minimum height=2.4cm,
              font=\tiny, align=center},
  lbl/.style={font=\scriptsize, align=center},
  arr/.style={-{Stealth[length=5pt]}, thick, blue!50}
]

%% Tabla original
\node[lbl] at (1.3, 1.7) {\textbf{Tabla original (25 M filas)}};
\node[col] (c1) at (0,    0) {cod\_mpio\\[2pt]05001\\05001\\05002\\\ldots};
\node[col] (c2) at (1.4,  0) {energia\\[2pt]1\\2\\1\\\ldots};
\node[col] (c3) at (2.8,  0) {paredes\\[2pt]2\\5\\1\\\ldots};

%% Flecha
\draw[arr] (3.6, 0) -- (5.0, 0) node[midway, above, font=\scriptsize]{VertiPaq};

%% Resultado comprimido
\node[lbl] at (7.5, 1.7) {\textbf{En memoria: diccionario + \'indices}};
\node[draw=orange!50, fill=orange!8, rounded corners, font=\tiny,
      inner sep=8pt, align=left] at (6.4, 0) {
  Diccionario ``energia'':\\
  \quad 0 $\to$ tiene servicio\\
  \quad 1 $\to$ no tiene servicio\\[4pt]
  25 M \'indices: 1 bit c/u\\
  $\approx$ \textbf{3 MB total}
};

\end{tikzpicture}
\end{center}

\medskip
{\footnotesize
\textbf{Regla clave de VertiPaq:}
cuanto \emph{menos} valores \'unicos tiene una columna,
m\'as peque\~na queda despu\'es de compresi\'on.\\
Una columna con 2 valores en 25 M filas ocupa $\approx$ 3 MB.
Una columna con 15 valores en \textbf{15 filas} ocupa $\approx$ \textbf{1 KB}.
}

\end{frame}

% =============================================================
\begin{frame}{?`Cu\'anto pesa una tabla de par\'ametro?}

\begin{columns}[T]

\column{0.50\textwidth}

\begin{block}{fct\_Sisben (tabla de hechos)}
\begin{tabular}{lr}
\toprule
Propiedad & Valor \\
\midrule
Filas     & 25{,}000{,}000 \\
Columnas  & 135 \\
Tama\~no comprimido & $\approx$ 1.5 GB \\
\bottomrule
\end{tabular}
\end{block}

\bigskip

\begin{block}{Dim\_IPM\_Dimension (tabla de par\'ametro)}
\begin{tabular}{lr}
\toprule
Propiedad & Valor \\
\midrule
Filas     & \textbf{15} \\
Columnas  & \textbf{2} \\
Tama\~no comprimido & $\approx$ \textbf{1 KB} \\
\bottomrule
\end{tabular}
\end{block}

\column{0.48\textwidth}
\begin{center}
\begin{tikzpicture}
  %% Barra grande = fct_Sisben
  \fill[blue!35] (0,0) rectangle (3.4,3.4);
  \node[white, font=\small\bfseries, align=center] at (1.7,1.7)
        {fct\_Sisben\\1.5 GB};

  %% Linea = tabla parametro (invisible en escala)
  \draw[orange!80!black, very thick] (0, 3.45) -- (3.4, 3.45);
  \node[font=\tiny, orange!70!black, align=left] at (1.7, 3.75)
        {Dim\_IPM\_Dimension $\approx$ 1 KB\\(invisible en esta escala)};

  %% Eje
  \draw[<->, gray, thin] (3.6,0) -- (3.6,3.4)
        node[midway, right, font=\tiny, gray]{1.5 GB};
\end{tikzpicture}
\end{center}

\bigskip
{\footnotesize
La tabla de par\'ametro es \textbf{1.500.000 veces\\m\'as peque\~na}.
Su impacto en RAM\\es literalmente \textbf{cero}.
}
\end{columns}

\end{frame}

% =============================================================
\begin{frame}{El problema real: mostrar 15 indicadores en una sola matriz}

\textbf{Situaci\'on:} 125 municipios, 15 indicadores IPM.
El usuario quiere ver la tasa de \emph{cualquiera} de los 15 por municipio.

\bigskip
\begin{columns}[T]

\column{0.48\textwidth}
\begin{exampleblock}{Opci\'on A --- 15 columnas fijas}
\begin{center}
\begin{tikzpicture}[font=\tiny]
  \fill[gray!15] (0,0) rectangle (4.2,1.4);
  \foreach \x in {0.3,0.6,...,4.0}{
    \draw[blue!30] (\x,0)--(\x,1.3);
    \node[rotate=90, blue!60!black] at (\x+0.12,0.65) {I\the\numexpr\x*10/3\relax};
  }
  \node[red!70!black, font=\tiny] at (2.1,-0.35)
        {15 columnas fijas, ilegible};
\end{tikzpicture}
\end{center}
\textcolor{red!70!black}{\textbf{Problema:}} el usuario no puede elegir; ve todo a la vez.
\end{exampleblock}

\column{0.48\textwidth}
\begin{alertblock}{Opci\'on B --- 15 p\'aginas}
\begin{itemize}\footnotesize
  \item P\'ag 1: municipios vs.\ bajo logro educativo
  \item P\'ag 2: municipios vs.\ analfabetismo
  \item \ldots
  \item P\'ag 15: municipios vs.\ hacinamiento
\end{itemize}
\textcolor{red!70!black}{\textbf{Problema:}} si la f\'ormula cambia,
hay que editar 15 visuals.
Si el t\'itulo cambia, 15 ediciones.
Esto se llama \emph{duplicaci\'on de l\'ogica}: el peor error en dise\~no BI.
\end{alertblock}

\end{columns}

\end{frame}

% =============================================================
\begin{frame}[fragile]{La soluci\'on: tabla de par\'ametro + SWITCH}

\begin{columns}[T]

\column{0.42\textwidth}

\textbf{Paso 1 --- La tabla de par\'ametro}\\
Solo guarda los \emph{nombres} de las opciones:

\bigskip
\begin{center}
\begin{tikzpicture}[font=\scriptsize]
\node[draw=orange!60, fill=orange!8, rounded corners,
      inner sep=8pt, align=left] {
  \textbf{Dim\_IPM\_Dimension}\\[2pt]
  \rule{4cm}{0.4pt}\\[2pt]
  \textit{Dimension\_IPM}\\[2pt]
  \rule{4cm}{0.4pt}\\[1pt]
  Bajo logro educativo\\
  Analfabetismo\\
  Inasistencia escolar\\
  Rezago escolar\\
  \ldots\ (15 filas en total)
};
\end{tikzpicture}
\end{center}

\textbf{Paso 2 ---} El usuario elige \\ ``Analfabetismo'' en el slicer.

\column{0.55\textwidth}

\textbf{Paso 3 --- La medida DAX captura la elecci\'on:}

\bigskip
\begin{tikzpicture}
\node[draw=blue!30, fill=blue!5, rounded corners,
      inner sep=8pt, font=\tiny\ttfamily,
      text width=6.4cm, align=left] {
  Pct\_Privacion\_Seleccionada =\\[2pt]
  SWITCH(\\
  \quad SELECTEDVALUE(\\
  \quad\quad Dim\_IPM\_Dimension[Dimension\_IPM]),\\[2pt]
  \quad "Bajo logro educativo",\\
  \quad\quad [Pct\_I1\_BajoLogroEducativo],\\
  \quad "Analfabetismo",\\
  \quad\quad \textbf{[Pct\_I2\_Analfabetismo],}\\
  \quad "Inasistencia escolar",\\
  \quad\quad [Pct\_I3\_InasistenciaEscolar],\\
  \quad ...\\
  )
};
\end{tikzpicture}

\bigskip
{\footnotesize $\Rightarrow$ La matriz muestra \textbf{un solo valor}
por municipio, el que corresponde a la elecci\'on del usuario.}

\end{columns}

\end{frame}

% =============================================================
\begin{frame}{Flujo completo: del clic del usuario al valor en pantalla}

\centering
\begin{tikzpicture}[
  dato/.style={rectangle, draw=orange!60, fill=orange!10, rounded corners=4pt,
               minimum height=1cm, minimum width=3cm,
               font=\small, align=center},
  calc/.style={rectangle, draw=blue!50, fill=blue!8, rounded corners=4pt,
               minimum height=1cm, minimum width=3.2cm,
               font=\small, align=center},
  viz/.style={rectangle, draw=green!50!black, fill=green!8, rounded corners=4pt,
              minimum height=1cm, minimum width=3cm,
              font=\small, align=center},
  arr/.style={-{Stealth[length=5pt]}, thick, gray!70}
]

\node[dato] (param) at (0,    0)
      {Slicer\\{\tiny Dim\_IPM\_Dimension}\\{\tiny 15 filas}};

\node[calc] (sel)   at (5.0,  0)
      {\texttt{SELECTEDVALUE()}\\{\tiny Lee ``Analfabetismo''}};

\node[calc] (sw)    at (5.0, -2)
      {\texttt{SWITCH()}\\{\tiny Mapea a [Pct\_I2]}};

\node[calc] (meas)  at (10.0, -2)
      {[Pct\_I2\_Analfabetismo]\\{\tiny Opera sobre fct\_Sisben}\\{\tiny 25 M filas}};

\node[viz]  (vis)   at (10.0,  0)
      {Matriz del tablero\\{\tiny Municipio | \% Analfab.}\\{\tiny Medell\'in | 4.2\%}};

%% Flechas
\draw[arr] (param) -- node[above, font=\tiny]{contexto de filtro} (sel);
\draw[arr] (sel)   -- (sw);
\draw[arr] (sw)    -- node[above, font=\tiny]{llama a la medida} (meas);
\draw[arr] (meas)  -- (vis);

%% Separador "sin relacion"
\draw[dashed, orange!60, thick] (2.3, 0.7) -- (2.3, -2.7);
\node[font=\tiny, orange!70!black, align=center] at (2.3, -3.1)
      {sin relaci\'on\\f\'isica entre\\tablas};

\end{tikzpicture}

\bigskip
{\footnotesize
\textbf{Punto clave:} la tabla de par\'ametro \textbf{nunca se une} a fct\_Sisben.
Funciona como un ``control remoto'': la medida DAX lee la selecci\'on
y decide qu\'e c\'alculo ejecutar sobre los 25 M de registros.\\
\src{SQLBI.com, ``Disconnected tables and SELECTEDVALUE'', Alberto Ferrari, 2020}
}

\end{frame}

% =============================================================
\begin{frame}{?`Por qu\'e funciona sin relaci\'on f\'isica?}

\begin{columns}[T]

\column{0.50\textwidth}

En Power BI las relaciones \textbf{propagan filtros} entre tablas.
Pero \texttt{SELECTEDVALUE()} \textbf{no necesita propagaci\'on}:
lee directamente el contexto de filtro del visual.

\bigskip
\begin{block}{An\'alogo: TV y control remoto}
\begin{center}
\begin{tikzpicture}[
  box/.style={draw, rounded corners, fill=gray!8,
              inner sep=6pt, minimum width=2.5cm,
              minimum height=0.8cm, font=\scriptsize, align=center},
  arr/.style={-{Stealth[length=4pt]}, thick}
]
\node[box] (tv)    at (3,  0)   {Pantalla\\(visual Power BI)};
\node[box] (ctrl)  at (0, -1.5) {Control remoto\\(slicer / tabla de param)};
\node[box] (datos) at (6, -1.5) {Cable de datos\\(fct\_Sisben)};

\draw[arr]          (ctrl)  -- node[left,  font=\tiny]{se\~nal} (tv);
\draw[arr]          (datos) -- node[right, font=\tiny]{datos}   (tv);
\draw[dashed, gray] (ctrl)  -- node[below, font=\tiny]{sin cable directo} (datos);
\end{tikzpicture}
\end{center}
El control cambia lo que muestra la pantalla,\\
sin estar conectado directamente al cable.
\end{block}

\column{0.47\textwidth}

\begin{alertblock}{?`Qu\'e es el contexto de filtro?}
Cuando el slicer selecciona ``Analfabetismo'',
Power BI establece un \textbf{contexto de filtro}
sobre \texttt{Dim\_IPM\_Dimension}.

\bigskip
\texttt{SELECTEDVALUE()} lee ese contexto:
\begin{itemize}\footnotesize
  \item Si hay \textbf{exactamente 1 valor} seleccionado $\to$ lo devuelve
  \item Si hay 0 o m\'as de 1 $\to$ devuelve BLANK
\end{itemize}

\bigskip
\textbf{No necesita relaci\'on} porque el contexto
de filtro no requiere relaci\'on f\'isica para
\emph{leerse}; solo para \emph{propagarse} a otras tablas.
\end{alertblock}

\bigskip
\src{Russo \& Ferrari, \textit{The Definitive Guide to DAX},\\Cap.~5: Filter Context, p.~142--178, 2019}

\end{columns}

\end{frame}

% =============================================================
\begin{frame}{Comparativa: tres formas de resolver el mismo problema}

\begin{center}
\begin{tabular}{lp{3cm}p{3cm}p{3.4cm}}
\toprule
 & \textbf{15 columnas fijas} & \textbf{15 p\'aginas} & \textbf{Tabla de par\'ametro} \\
\midrule
El usuario elige qu\'e ver
  & \textcolor{red!70!black}{No}
  & \textcolor{orange}{Solo por p\'agina}
  & \textcolor{green!50!black}{\textbf{S\'i, en tiempo real}} \\
\addlinespace
Visuals necesarios
  & 1 (ilegible)
  & 15
  & \textbf{1} \\
\addlinespace
Si cambia la f\'ormula
  & Editar 1 medida
  & \textcolor{red!70!black}{Editar 15 visuals}
  & \textbf{Editar 1 medida} \\
\addlinespace
Peso en RAM extra
  & 0
  & 0
  & \textbf{$\sim$1 KB} \\
\addlinespace
Mantenimiento
  & Bajo
  & \textcolor{red!70!black}{Muy alto}
  & \textbf{Bajo} \\
\addlinespace
Interactividad
  & Ninguna
  & Baja
  & \textbf{Alta} \\
\bottomrule
\end{tabular}
\end{center}

\bigskip
\begin{exampleblock}{Conclusi\'on}
La tabla de par\'ametro a\~nade \textbf{$\approx$ 1 KB de peso} y
\textbf{interactividad total}.
Es el patr\'on recomendado por Microsoft y SQLBI para matrices din\'amicas.\\[4pt]
\src{Microsoft, ``Field parameters in Power BI Desktop'', docs.microsoft.com, 2022}
\end{exampleblock}

\end{frame}

% =============================================================
\begin{frame}{Resumen}

\begin{enumerate}\setlength\itemsep{0.6em}

  \item \textbf{VertiPaq comprime por columna, no por fila.}
    Una tabla de 15 filas ocupa $\approx$ 1 KB en memoria,
    sin importar que la tabla de hechos tenga 25 millones de filas.

  \item \textbf{La tabla de par\'ametro no necesita relaci\'on.}
    Su funci\'on es solo proveer opciones al slicer.
    La medida DAX lee la elecci\'on del usuario con
    \texttt{SELECTEDVALUE()}, sin propagaci\'on de filtros.

  \item \textbf{Sin ella, las alternativas son peores.}
    15 columnas fijas: ilegible.
    15 p\'aginas: duplicaci\'on de l\'ogica, mantenimiento insostenible.

  \item \textbf{Es el est\'andar de la industria.}
    Microsoft lo document\'o como \textit{Field Parameters} en 2022.
    SQLBI lo recomienda como pr\'actica \'optima desde 2015.

\end{enumerate}

\bigskip
\begin{center}
\src{Fuentes: Russo \& Ferrari (2019) $\cdot$
SQLBI.com, Ferrari (2015, 2020) $\cdot$
Microsoft Docs (2022)}
\end{center}

\end{frame}

\end{document}
